\documentclass{article} % amsmart?

\usepackage{amsthm}
\usepackage{amsmath}
\usepackage{graphicx}
\usepackage{amssymb}

\newtheorem{proposition}{Proposition}
\newtheorem{lemma}{Lemma}

\title{MATH290}
\author{Caleb Hessing}
\date{Jan 1, 1970}

\begin{document}

\maketitle

\begin{proposition}
    For any integer $k$, $k^3 - k$ is a multiple of $6$.
\end{proposition}

\begin{proof}\nobreakspace

    Firstly, note that the sum of two multiples of $6$ is itself a multiple of $6$. An integer
    is a multiple of $6$ if it can be expressed as the product of $6$ and some other integer. 
    Consider a multiple of six $n = 6i$ and another $m = 6j$. Observe,

    \[
    \begin{aligned}
    n + m &= 6i + 6j \\
    &= 6(i + j)
    \end{aligned}
    \]

    Relying on the property that summing integers yields an integer, we've shown that $n + m$
    can be expressed as a product of $6$ and some other integer $i + j$. That is, $n+ m$ is
    a multiple of $6$ for any multiples of six $n$ and $m$.

    Secondly, observe what happens if you substitute $-k$ for $k$:

    \[
    \begin{aligned}
        & (-k)^3 - (-k) \\
        =& -(k^3) + k \\
        =& -(k^3 - k)
    \end{aligned}
    \]

    Thus, if $k^3 - k$ can be expressed as $6n$ for some integers $k$ and $n$, then $(-k)^3 - (-k)$
    can be expressed as $-6n$. Therefore, if I can show that $k^3 - k$ is a multiple of $6$ for
    all integers $k \ge 0$, then $k^3 - k$ is a multiple of $6$ for all integers $k < 0$. With
    this in mind, I'll only focus on integers $k \ge 0$ for the rest of the proof.

    Let $f_k$ be the sequence $k^3 - k$ for integers $k \ge 0$. It turns out that rather than
    proving each term in this sequence is a multiple of $6$, it's easier to prove that
    \textit{the difference between consecutive terms} is a multiple of $6$. (Or at least
    \textit{I} think it is.) Let's let $g_k$ be the difference between any $f_{k+1}$ and 
    $f_{k}$. That is,

    \[
    \begin{aligned}
        g_k &= f_{k+1} - f_{k} \\
        f_{k+1} &= f_k + g_k
    \end{aligned}
    \]

    If I can show that $g_k$ is a multiple of $6$ for all integers $k$, then so long as the
    sequence starts with a multiple of $6$ all following terms will be multiples of $6$. 


\end{proof}

\end{document}